% TOP_PROBLEM: {"id": 3186, "title": "Forcing a $K_6$-minor", "category_id": 3, "category": {"id": 3, "name": "graph_theory", "display_name": "Graph Theory", "description": "Problems involving graphs, networks, and their properties.", "slug": "graph-theory", "order_index": 3, "created_at": "2026-07-31T15:26:25.671Z"}, "status": "open", "problem_number": "OPG-37379", "set_id": 12}
% FIRST_POSTED: 2026-10-01
% ENTRY_KIND: statement_only
% UnsolvedMath ID 3186; source catalogue code OPG-37379
% !TeX program = lualatex
% Definition and sources only; this file is not an LLM solution attempt.
\documentclass[11pt,a4paper]{article}
\usepackage[margin=25mm]{geometry}
\usepackage{fontspec}
\setmainfont{lmroman10-regular.otf}[BoldFont=lmroman10-bold.otf,
 ItalicFont=lmroman10-italic.otf,BoldItalicFont=lmroman10-bolditalic.otf]
\usepackage{amsmath,amssymb,mathtools}
\usepackage{unicode-math}
\setmathfont{latinmodern-math.otf}
\AtBeginDocument{\let\setminus\smallsetminus}
\usepackage{xurl,hyperref}
\hypersetup{colorlinks=true,urlcolor=blue}
\setcounter{secnumdepth}{0}
\setlength{\parindent}{0pt}
\setlength{\parskip}{5pt}
\setlength{\emergencystretch}{3em}
\begin{document}
\section*{Forcing a $K_6$-minor}\label{3186}
{\small UnsolvedMath ID 3186 · OPG-37379 · Graph Theory · OpenGarden}
\subsection{Definitions and mathematical statement}
Let $G=(V,E)$ be a finite simple undirected graph. Its minimum degree is $\delta(G)=\min_{v\in V}\deg(v)$. It contains a $K_6$ minor if there are six pairwise disjoint nonempty subsets $B_1,\ldots,B_6\subseteq V$ such that each induced subgraph $G[B_i]$ is connected and, for every $i\ne j$, at least one edge joins $B_i$ to $B_j$. The branch sets may have arbitrary sizes. The edges witnessing different pairs need not meet distinct vertices within a branch set.

The source presents two assertions:
\begin{enumerate}
\item Every finite simple graph of minimum degree at least seven contains a $K_6$ minor.
\item Every finite simple seven-vertex-connected graph contains a $K_6$ minor.
\end{enumerate}
Seven-vertex-connected means having at least eight vertices and remaining connected after deleting any set of at most six vertices. The first assertion implies the second, since this connectivity condition forces minimum degree at least seven; the converse implication is not supplied by the definitions.

The minimum-degree hypothesis concerns every vertex individually, and is not an average-degree assumption. Likewise, a minor need not occur as a complete subgraph or as a subdivision with internally disjoint paths. A disproof of the first assertion must meet the degree bound but may have low connectivity; a disproof of the second would necessarily disprove both. Each assertion is quantified over graphs of all finite orders.
\subsection{Short English statement}
Must minimum degree seven force a complete six-vertex minor, and does seven-vertex-connectivity suffice?
\subsection{Sources}
\begin{itemize}
\item[S1] ULAM.AI and the UnsolvedMath Contributors, supplied problems.json, record 3186 (OPG-37379), OpenGarden. Catalogue identity and source transcription; CC BY 4.0.
\url{https://huggingface.co/datasets/ulamai/UnsolvedMath}.
\item[S2] Open Problem Garden, Forcing a $K_6$-minor. Original problem and discussion.
\url{http://www.openproblemgarden.org/op/forcing_a_k_6_minor}.
\end{itemize}
\end{document}
